\chapter{Elektronische effecten en reactieve intermediairen}\label{ch:b1:electronic-effects}

Trifluorethaanzuur, \ce{CF3COOH}, is als zuur ongeveer tienduizend keer
sterker dan ethaanzuur, \ce{CH3COOH}, hoewel de zure O--H-binding drie
bindingen van de fluoratomen verwijderd is. Fenol is een miljoen keer
zuurder dan ethanol, en een ammoniakmolecuul dat aan een benzeenring
gebonden is --- aniline --- is een miljoen keer minder basisch dan een
dat aan een methylgroep gebonden is. De organische chemie zit vol met
zulke contrasten, en twee effecten verklaren de meeste ervan: de trek van
elektronegatieve atomen langs de $\sigma$-bindingen, en de spreiding van
elektronenparen over $\pi$-systemen. Dezelfde twee effecten bepalen welke
bindingen breken, welke intermediairen ontstaan en waar reagentia
aanvallen: ze zijn de grammatica van de \omterm{def:b1:elementary-steps:elementary-step}{reactiemechanismen} van de
volgende hoofdstukken.

\begin{recall}
\omterm{def:b1:periodicity:electronegativity}{Elektronegativiteit} en de trends ervan (\cref{ch:b1:periodicity});
lewisstructuren, \omterm{def:b1:lewis-resonance-vsepr:formal-charge}{formele ladingen} en resonantie
(\cref{ch:b1:lewis-resonance-vsepr}); $\mathrm{p}K_a$ en de sterkte van
zuren en basen (\cref{ch:b1:predominant-reaction}). Het schoolboek (jaar
12) tekende \omterm{def:b1:electronic-effects:curly-arrow}{gebogen pijlen} van een \omterm{def:b1:electronic-effects:nucleophile}{nucleofiel} naar een \omterm{def:b1:electronic-effects:nucleophile}{elektrofiel}.
\end{recall}

\section{Inductieve effecten}

\begin{definition}[Inductief effect]\label{def:b1:electronic-effects:inductive}
Het \emph{inductieve effect}\index{inductief effect} is de verschuiving
van elektronendichtheid langs $\sigma$-bindingen naar een elektronegatief
atoom of een elektronegatieve groep. Een elektronenzuigende groep (F, Cl,
O, \ce{NO2}, \ce{CF3}) heeft een $-I$-effect; alkylgroepen en negatief
geladen atomen stuwen elektronen, een $+I$-effect. Het effect verzwakt
snel met het aantal tussenliggende bindingen.
\end{definition}

\begin{proposition}[Inductief effect en zuurgraad]\label{prop:b1:electronic-effects:inductive-pka}
Een elektronenzuigende groep in de buurt van de zure plaats van een
carbonzuur stabiliseert het carboxylaation door zijn negatieve lading te
spreiden, en verlaagt de $\mathrm{p}K_a$: hoe meer zuigende groepen en hoe
dichterbij, hoe sterker het zuur.
\end{proposition}

\begin{proof}
De $\mathrm{p}K_a$ meet de ligging van $\mathrm{RCOOH} \rightleftharpoons
\mathrm{RCOO^-} + \ce{H+}$. Een groep die elektronendichtheid van het
carboxylaat wegtrekt, spreidt de lading ervan over meer atomen, wat de
energie ervan sterker verlaagt dan die van het neutrale zuur; het
evenwicht verschuift naar rechts en $K_a$ neemt toe. De gemeten reeks
bevestigt dat.
\end{proof}

\begin{omfigure}
\begin{tikzpicture}
\begin{axis}[
  xbar, width=0.78\linewidth, height=5.2cm, xmin=0, xmax=5.4,
  symbolic y coords={TFA,TCA,DCA,MCA,AA}, ytick=data,
  yticklabels={\ce{CF3COOH}, \ce{CCl3COOH}, \ce{CHCl2COOH}, \ce{CH2ClCOOH}, \ce{CH3COOH}},
  xlabel={$\mathrm{p}K_a$ bij \qty{25}{\celsius}}, bar width=9pt,
  nodes near coords, nodes near coords align={horizontal},
  every node near coord/.append style={font=\scriptsize},
  tick label style={font=\small}, label style={font=\small}, enlarge y limits=0.15]
\addplot[fill=omDef!45, draw=omDef] coordinates {(0.52,TFA) (0.51,TCA) (1.35,DCA) (2.87,MCA) (4.76,AA)};
\end{axis}
\end{tikzpicture}

{\small De gehalogeneerde ethaanzuren. Elk chlooratoom verlaagt de
$\mathrm{p}K_a$; met drie is het zuur ongeveer tienduizend keer sterker
dan ethaanzuur. Trifluor- en trichloorethaanzuur zijn binnen de
onzekerheid van metingen aan zulke \omterm{def:b1:predominant-reaction:strong-acid}{sterke zuren} even sterk.}
% ledger: pka:ethanoic, pka:chloroethanoic, pka:dichloroethanoic, pka:trichloroethanoic, pka:trifluoroethanoic
\end{omfigure}

Alkylgroepen werken een beetje in de andere richting: propaanzuur
($\mathrm{p}K_a
= 4.89$) is iets zwakker dan ethaanzuur (4.76), en methaanzuur (3.74),
zonder alkylgroep, is sterker dan beide.
% ledger: pka:propanoic, pka:methanoic


\section{Mesomere effecten}

Wanneer een atoom met een \omterm{def:b1:lewis-resonance-vsepr:lewis}{vrij elektronenpaar}, of met een lading, aan een
$\pi$-systeem gebonden is, kunnen zijn elektronen met dat systeem gedeeld
worden; \omterm{def:b1:lewis-resonance-vsepr:resonance}{resonantiestructuren} (\cref{ch:b1:lewis-resonance-vsepr})
beschrijven die deling.

\begin{definition}[Mesomeer effect, donor- en acceptorgroepen]\label{def:b1:electronic-effects:mesomeric}
Het \emph{mesomere effect}\index{mesomeer effect} is de verschuiving van
elektronendichtheid via een geconjugeerd $\pi$-systeem, weergegeven met
\omterm{def:b1:lewis-resonance-vsepr:resonance}{resonantiestructuren}. Een \emph{donorgroep}\index{donorgroep} geeft een
\omterm{def:b1:lewis-resonance-vsepr:lewis}{vrij elektronenpaar} aan het systeem, een $+M$-effect (\ce{-OH},
\ce{-O^-}, \ce{-NH2}, \ce{-OR}); een
\emph{acceptorgroep}\index{acceptorgroep} trekt er elektronen uit via een
$\pi$-binding met een elektronegatief atoom, een $-M$-effect (\ce{-NO2},
\ce{-C=O}, \ce{-C#N}).
\end{definition}

\begin{omfigure}
\setchemfig{atom sep=2em}
\schemestart
\chemfig{CH_3-C(=[1]O)-[7]O^{-}}
\arrow{<->}
\chemfig{CH_3-C(-[1]O^{-})=[7]O}
\schemestop
\hspace{2.5em}
\schemestart
\chemfig{*6(=-=(-O^{-})-=-)}
\arrow{<->}
\chemfig{*6((=O)-=-\chemabove{}{\scriptstyle\ominus}-=-)}
\schemestop

\medskip
\schemestart
\chemfig{H_3C-C(=[1]O)-[7]NH_2}
\arrow{<->}
\chemfig{H_3C-C(-[1]O^{-})=[7]NH_2^{+}}
\schemestop

{\small \omterm{def:b1:lewis-resonance-vsepr:resonance}{Resonantiestructuren} van het ethanoaation (de lading verdeeld over
twee gelijkwaardige zuurstofatomen), van het fenoxide-ion (de lading
gespreid in de ring; één van de drie ringstructuren getekend), en van een
amide (het \omterm{def:b1:lewis-resonance-vsepr:lewis}{vrije paar} van stikstof gedeeld met de C=O, wat amiden tot
\omterm{def:b1:predominant-reaction:strong-acid}{zwakke basen} maakt).}
\end{omfigure}

\begin{proposition}[Mesomeer effect en zuurgraad]\label{prop:b1:electronic-effects:mesomeric-pka}
Een zuur waarvan de geconjugeerde base haar lading door resonantie
spreidt, is sterker dan een zuur waarvan de base dat niet kan:
carbonzuren ($\mathrm{p}K_a$ ongeveer 4--5) zijn veel sterker dan
alcoholen (ongeveer 16), en fenol (10.0) ligt ertussen. Een \omterm{def:b1:electronic-effects:mesomeric}{acceptorgroep}
op de \textit{ortho}- of \textit{para}-positie ten opzichte van de OH van
een fenol versterkt het meer dan op de \textit{meta}-positie.
\end{proposition}

\begin{proof}
In een alkoxide zit de lading op één zuurstofatoom; in een carboxylaat
wordt ze gedeeld door twee zuurstofatomen; in een fenoxide door het
zuurstofatoom en drie ringkoolstofatomen, die minder elektronegatief zijn
dan zuurstof, een kleinere winst. Een \ce{NO2}-groep op het
\textit{ortho}- of \textit{para}-koolstofatoom kan de lading zelf opnemen
via een \omterm{def:b1:lewis-resonance-vsepr:resonance}{resonantiestructuur} met \ce{C=N+(-O^-)-O^-}; bij \textit{meta}
brengt geen enkele \omterm{def:b1:lewis-resonance-vsepr:resonance}{resonantiestructuur} de lading op het koolstofatoom dat
\ce{NO2} draagt, en blijft alleen het \omterm{def:b1:electronic-effects:inductive}{inductieve effect} over. De gemeten
waarden: 2-nitrofenol 7.23, 4-nitrofenol 7.16, 3-nitrofenol 8.36, tegen
9.99 voor fenol en 15.9 voor ethanol.
\end{proof}
% ledger: pka:ethanol, pka:phenol, pka:2-nitrophenol, pka:3-nitrophenol, pka:4-nitrophenol

Dezelfde redenering geldt voor basen. Methylamine ($\mathrm{p}K_a$ van
zijn ammonium 10.66) is een sterkere base dan ammoniak (9.25,
\cref{ch:b1:predominant-reaction}), omdat de methylgroep elektronen naar
stikstof stuwt; in aniline wordt het \omterm{def:b1:lewis-resonance-vsepr:lewis}{vrije paar} met de ring gedeeld en is
het ammonium veel zuurder (4.60). Pyridine (5.23) behoudt zijn \omterm{def:b1:lewis-resonance-vsepr:lewis}{vrije
paar}, maar op een stikstofatoom in een vlakke ring met meer
\textit{s}-karakter, minder beschikbaar dan in een amine.
% ledger: pka:methylammonium, pka:anilinium, pka:pyridinium, dfg:NH4+_ao


\section{Een binding breken: reactieve intermediairen}

\begin{definition}[Homolyse en heterolyse]\label{def:b1:electronic-effects:cleavage}
Een \omterm{def:b1:lewis-resonance-vsepr:lewis}{covalente binding} breekt door \emph{homolyse}\index{homolyse} wanneer
elk atoom één van de twee bindingselektronen houdt, wat twee \omterm{def:b1:electronic-effects:intermediates}{radicalen}
geeft, en door \emph{heterolyse}\index{heterolyse} wanneer één atoom
beide houdt, wat een kation en een anion geeft.
\end{definition}

\begin{definition}[Reactieve intermediairen]\label{def:b1:electronic-effects:intermediates}
Een \emph{reactief intermediair}\index{reactief intermediair} is een
deeltje dat in één stap van een mechanisme gevormd en in een latere stap
verbruikt wordt, te reactief om zich op te hopen
(\cref{ch:b1:elementary-steps}). Een
\emph{carbokation}\index{carbokation} heeft een koolstofatoom met drie
bindingen en een positieve lading (zes \omterm{def:b1:quantum-numbers:configuration}{valentie-elektronen}); een
\emph{carbanion}\index{carbanion} een koolstofatoom met drie bindingen,
een \omterm{def:b1:lewis-resonance-vsepr:lewis}{vrij elektronenpaar} en een negatieve lading; een
\emph{radicaal}\index{radicaal} een atoom met een \omterm{def:b1:quantum-numbers:unpaired}{ongepaard elektron}.
\end{definition}

\begin{definition}[Klasse van een koolstofatoom]\label{def:b1:electronic-effects:carbon-class}
Een koolstofatoom dat aan één ander koolstofatoom gebonden is, is een
\emph{primair koolstofatoom}\index{primair koolstofatoom}, aan twee een
\emph{secundair koolstofatoom}\index{secundair koolstofatoom}, aan drie
een \emph{tertiair koolstofatoom}\index{tertiair koolstofatoom}. Een
\omterm{def:b1:electronic-effects:intermediates}{carbokation}, een \omterm{def:b1:electronic-effects:intermediates}{radicaal} of een halogeenalkaan krijgt de klasse van het
betrokken koolstofatoom.
\end{definition}

\begin{omfigure}
\begin{tikzpicture}[font=\small]
  % carbocation: planar, empty p orbital
  \begin{scope}
    \fill[omProp!15, draw=omProp] (0,0.15) ellipse (0.18 and 0.75);
    \fill[omProp!15, draw=omProp] (0,-0.15) ellipse (0.18 and 0.75);
    \draw[thick] (0,0) -- (-1.0,-0.25) node[left] {R};
    \draw[thick] (0,0) -- (0.95,-0.35) node[right] {R};
    \draw[thick] (0,0) -- (0.35,0.45) node[above right] {R};
    \node[fill=white, inner sep=1pt] at (0,0) {\ce{C+}};
    \node at (0,-1.75) {carbokation: vlak,};
    \node at (0,-2.15) {lege \textit{p}-orbitaal};
  \end{scope}
  % radical: nearly planar, one electron
  \begin{scope}[xshift=4.6cm]
    \draw[omMeth] (0,0.15) ellipse (0.18 and 0.75);
    \draw[omMeth] (0,-0.15) ellipse (0.18 and 0.75);
    \fill (0,0.55) circle (0.05);
    \draw[thick] (0,0) -- (-1.0,-0.25) node[left] {R};
    \draw[thick] (0,0) -- (0.95,-0.35) node[right] {R};
    \draw[thick] (0,0) -- (0.35,0.45) node[above right] {R};
    \node[fill=white, inner sep=1pt] at (0,0) {C};
    \node at (0,-1.75) {radicaal: bijna vlak,};
    \node at (0,-2.15) {één elektron};
  \end{scope}
  % carbanion: pyramidal with a lone pair
  \begin{scope}[xshift=9.2cm]
    \draw[thick] (0,0) -- (-0.95,-0.55) node[left] {R};
    \draw[thick] (0,0) -- (0.95,-0.55) node[right] {R};
    \draw[thick] (0,0) -- (0.2,-0.9) node[below] {R};
    \fill[omDef!15, draw=omDef] (0,0.55) ellipse (0.22 and 0.45);
    \fill (-0.07,0.65) circle (0.045); \fill (0.07,0.65) circle (0.045);
    \node[fill=white, inner sep=1pt] at (0,0) {\ce{C-}};
    \node at (0,-1.75) {carbanion: piramidaal,};
    \node at (0,-2.15) {vrij elektronenpaar};
  \end{scope}
\end{tikzpicture}

{\small Geometrieën van de drie \omterm{def:b1:electronic-effects:intermediates}{reactieve intermediairen} van koolstof (R:
H of een alkylgroep). Het vlakke \omterm{def:b1:electronic-effects:intermediates}{carbokation} kan aan beide zijden
aangevallen worden, een punt dat belangrijk is voor de stereochemie van
\cref{ch:b1:nucleophilic-substitution}.}
\end{omfigure}

\begin{proposition}[Stabiliteit van carbokationen]\label{prop:b1:electronic-effects:carbocation-order}
\omterm{def:b1:electronic-effects:intermediates}{Carbokationen} worden gestabiliseerd door groepen die hun elektronen
geven: tertiair $>$ secundair $>$ primair $>$ methyl; een \omterm{def:b1:electronic-effects:intermediates}{carbokation}
naast een dubbele binding (allylisch) of een benzeenring (benzylisch)
wordt bovendien door resonantie gestabiliseerd, en een \omterm{def:b1:electronic-effects:intermediates}{carbokation} naast
een atoom met een \omterm{def:b1:lewis-resonance-vsepr:lewis}{vrij elektronenpaar} (\ce{R-O-CH2+}) nog meer.
\end{proposition}

\begin{proof}
Elke alkylgroep stuwt elektronendichtheid naar het elektronenarme
koolstofatoom ($+I$-effect, en een verwante wisselwerking van zijn
C--H-bindingen met de lege orbitaal, die in het deel van jaar~2 beschreven
wordt). Een allylisch kation heeft twee \omterm{def:b1:lewis-resonance-vsepr:resonance}{resonantiestructuren} die de
lading op het ene of het andere eindkoolstofatoom leggen; een
benzylisch kation heeft er vier, drie in de ring; een \omterm{def:b1:lewis-resonance-vsepr:lewis}{vrij paar} van
zuurstof vormt een vierde binding naar het kationische koolstofatoom,
\ce{R-O+=CH2}, waarin elk atoom een octet heeft.
\end{proof}

\omterm{def:b1:electronic-effects:intermediates}{Carbanionen} volgen de omgekeerde volgorde (methyl $>$ primair $>$
secundair $>$ tertiair) en worden gestabiliseerd door \omterm{def:b1:electronic-effects:mesomeric}{acceptorgroepen};
\omterm{def:b1:electronic-effects:intermediates}{radicalen} volgen dezelfde volgorde als \omterm{def:b1:electronic-effects:intermediates}{carbokationen}, minder uitgesproken.

\section{Gebogen pijlen, nucleofielen en elektrofielen}

\begin{definition}[Gebogen pijlen]\label{def:b1:electronic-effects:curly-arrow}
In een mechanisme toont een \emph{gebogen pijl}\index{gebogen pijl} de
verplaatsing van een elektronenpaar: hij begint op een \omterm{def:b1:lewis-resonance-vsepr:lewis}{vrij paar} of een
binding en eindigt op een atoom (waar een \omterm{def:b1:lewis-resonance-vsepr:lewis}{vrij paar} of een nieuwe binding
ontstaat) of tussen twee atomen (een nieuwe binding). Een \emph{halve
pijl}\index{halve pijl} toont de verplaatsing van één enkel elektron, in
radicaalreacties.
\end{definition}

\begin{method}[Gebogen pijlen tekenen]\label{met:b1:electronic-effects:curly-arrows}
\begin{enumerate}
  \item Vertrek van elektronen: een \omterm{def:b1:lewis-resonance-vsepr:lewis}{vrij paar}, een $\pi$-binding of een
        $\sigma$-binding, nooit van een positieve lading of van een atoom
        zonder elektronen.
  \item Eindig waar het paar heen gaat: een nieuwe binding tussen de twee
        atomen, of een \omterm{def:b1:lewis-resonance-vsepr:lewis}{vrij paar} op het meest elektronegatieve atoom
        wanneer een binding breekt.
  \item Tel de ladingen na elke stap: het atoom dat een \omterm{def:b1:lewis-resonance-vsepr:lewis}{vrij paar} aan een
        nieuwe binding geeft, krijgt er $+1$ bij, het atoom dat een paar
        als \omterm{def:b1:lewis-resonance-vsepr:lewis}{vrij paar} ontvangt, $-1$; de totale lading blijft behouden.
  \item Overschrijd nooit een octet op C, N, O: wanneer er een paar
        aankomt, moet er een ander vertrekken.
\end{enumerate}
\end{method}

\begin{omfigure}
\setchemfig{atom sep=2.2em}
\schemestart
\chemfig{H_3C-C(-[2]H)(-[6]H)-[@{b}]@{br}Br}
\arrow{->[heterolyse]}[,1.5]
\chemfig{H_3C-C(-[2]H)(-[6]H)^{+}}\+\chemfig{Br^{-}}
\schemestop
\chemmove{\draw[->, shorten <=2pt, shorten >=2pt](b) .. controls +(90:6mm) and +(90:6mm) .. (br);}

\medskip
\schemestart
\chemfig{H@{o}O^{-}}\+\chemfig{@{h}H-@{oc}O-H}
\arrow{<=>}
\chemfig{H_2O}\+\chemfig{HO^{-}}
\schemestop
\chemmove{\draw[->, thick, shorten <=2pt, shorten >=2pt](o) .. controls +(-80:8mm) and +(-100:8mm) .. (h);}
\par\vspace{5mm}

{\small \omterm{def:b1:electronic-effects:curly-arrow}{Gebogen pijlen}. Boven: \omterm{def:b1:electronic-effects:cleavage}{heterolyse} van de C--Br-binding, waarbij
het paar naar broom gaat. Onder: een protonoverdracht van water naar
hydroxide, geschreven als de aanval van een \omterm{def:b1:lewis-resonance-vsepr:lewis}{vrij paar} van de base op het
waterstofatoom (de O--H-binding van water breekt dan, en haar paar blijft
op zuurstof; de tweede pijl wordt overgelaten aan
\cref{exo:b1:electronic-effects:4}).}
\end{omfigure}

\begin{definition}[Nucleofiel, elektrofiel, vertrekkende groep]\label{def:b1:electronic-effects:nucleophile}
Een \emph{nucleofiel}\index{nucleofiel} is een deeltje dat een
elektronenpaar geeft aan een ander atoom dan waterstof om een binding te
vormen (een \omterm{def:b1:lewis-resonance-vsepr:lewis-acid}{lewisbase}); een \emph{elektrofiel}\index{elektrofiel} is een
deeltje dat zo'n paar opneemt (een \omterm{def:b1:lewis-resonance-vsepr:lewis-acid}{lewiszuur}, vaak een koolstofatoom met
een positieve partiële lading). De
\emph{nucleofiliciteit}\index{nucleofiliciteit} van een deeltje is zijn
reactiviteit als nucleofiel, gemeten met \omterm{def:b1:rate-laws:order}{snelheidsconstanten}. Een
\emph{vertrekkende groep}\index{vertrekkende groep} is de groep die met
het bindende paar vertrekt wanneer een binding met koolstof heterolytisch
breekt.
\end{definition}

\begin{proposition}[Nucleofiliciteit en basiciteit]\label{prop:b1:electronic-effects:nucleophilicity-basicity}
Voor \omterm{def:b1:electronic-effects:nucleophile}{nucleofielen} die via hetzelfde atoom aanvallen, volgt de
\omterm{def:b1:electronic-effects:nucleophile}{nucleofiliciteit} de basiciteit: \ce{CH3O-} $>$ \ce{OH-} $>$ \ce{CH3COO-}
$>$ \ce{H2O}. Dat gaat niet op binnen een kolom van het periodiek systeem
in \omterm{def:b1:intermolecular-forces-solvents:solvent-classes}{protische oplosmiddelen}, waar grote, polariseerbare, zwak
gesolvateerde \omterm{def:b1:electronic-effects:nucleophile}{nucleofielen} sneller reageren hoewel ze zwakkere basen zijn
(\ce{I-} $>$ \ce{Br-} $>$ \ce{Cl-} $>$ \ce{F-}), en evenmin voor
omvangrijke basen, die slechte \omterm{def:b1:electronic-effects:nucleophile}{nucleofielen} zijn.
\end{proposition}

\begin{proof}
De basiciteit meet het evenwicht van de aanval op een proton, de
\omterm{def:b1:electronic-effects:nucleophile}{nucleofiliciteit} de snelheid van de aanval op een koolstofatoom; beide
nemen toe met de beschikbaarheid van het elektronenpaar, vandaar het
parallellisme binnen één atoomsoort. De snelheid hangt ook af van de
energie die nodig is om het oplosmiddel van het \omterm{def:b1:electronic-effects:nucleophile}{nucleofiel} te strippen en
van de vervorming van zijn elektronenwolk naar een overvol
koolstofatoom: kleine anionen zoals \ce{F-} worden sterk gebonden door
moleculen van een \omterm{def:b1:intermolecular-forces-solvents:solvent-classes}{protisch oplosmiddel}
(\cref{ch:b1:intermolecular-forces-solvents}), en omvangrijke groepen
hinderen de nadering van koolstof, maar niet die van een proton.
\end{proof}

Een goede \omterm{def:b1:electronic-effects:nucleophile}{vertrekkende groep} is een \omterm{def:b1:predominant-reaction:strong-acid}{zwakke base}, stabiel zodra ze het
paar draagt: \ce{I-}, \ce{Br-}, \ce{Cl-}, water, sulfonaten; \ce{OH-} en
\ce{RO-}, \omterm{def:b1:predominant-reaction:strong-acid}{sterke basen}, zijn slechte \omterm{def:b1:electronic-effects:nucleophile}{vertrekkende groepen} --- de reden
waarom alcoholen geactiveerd moeten worden vóór een substitutie
(\cref{ch:b1:alcohol-activation}).

\section{Oefeningen}

\begin{exercise}[$\star$]\label{exo:b1:electronic-effects:1}
Deel elk koolstofatoom van 2-methylbutaan en van
2-broom-2-methylpropaan in als primair, secundair, tertiair of quaternair
(gebonden aan vier koolstofatomen). Welk \omterm{def:b1:electronic-effects:intermediates}{carbokation} ontstaat het
gemakkelijkst uit elk bromide met formule \ce{C4H9Br}?
\end{exercise}

\begin{exercise}[$\star$]\label{exo:b1:electronic-effects:2}
Bereken uit de $\mathrm{p}K_a$-waarden de verhouding van de \omterm{def:b1:predominant-reaction:acidity-constant}{zuurconstanten}
van chloorethaanzuur en ethaanzuur, en van trifluorethaanzuur en
ethaanzuur.
\end{exercise}

\begin{exercise}[$\star$]\label{exo:b1:electronic-effects:3}
Teken de \omterm{def:b1:lewis-resonance-vsepr:resonance}{resonantiestructuren} van het ethanoaation, van het
4-nitrofenoxide-ion (met de lading op een zuurstofatoom van de
nitrogroep) en van het allylkation \ce{CH2=CH-CH2+}.
\end{exercise}

\begin{exercise}[$\star$]\label{exo:b1:electronic-effects:4}
Vervolledig de \omterm{def:b1:electronic-effects:curly-arrow}{gebogen pijlen} van de protonoverdracht van de figuur en
van \ce{NH3 + H3O+ -> NH4+ + H2O}. Controleer de ladingen.
\end{exercise}

\begin{exercise}[$\star\star$]\label{exo:b1:electronic-effects:5}
Rangschik naar toenemende zuurgraad: ethanol, fenol, ethaanzuur,
4-nitrofenol, trichloorethaanzuur, water ($\mathrm{p}K_a$ van het koppel
\ce{H2O}/\ce{OH-}, 14.0). Verantwoord elke stap met een inductief of
mesomeer argument.
\end{exercise}

\begin{exercise}[$\star\star$]\label{exo:b1:electronic-effects:6}
Rangschik naar toenemende basiciteit: aniline, methylamine, ammoniak,
pyridine. Verklaar de plaats van aniline met een \omterm{def:b1:lewis-resonance-vsepr:resonance}{resonantiestructuur}.
\end{exercise}

\begin{exercise}[$\star\star$]\label{exo:b1:electronic-effects:7}
Waarom is 3-nitrofenol zwakker dan 4-nitrofenol, maar toch sterker dan
fenol? Teken de \omterm{def:b1:lewis-resonance-vsepr:resonance}{resonantiestructuren} die beide uitspraken
rechtvaardigen.
\end{exercise}

\begin{exercise}[$\star\star$]\label{exo:b1:electronic-effects:8}
Rangschik de \omterm{def:b1:electronic-effects:intermediates}{carbokationen} \ce{CH3+}, \ce{(CH3)3C+}, \ce{CH2=CH-CH2+},
\ce{CH3CH2+}, \ce{C6H5CH2+}, \ce{CH3OCH2+}, en verantwoord.
\end{exercise}

\begin{exercise}[$\star\star$]\label{exo:b1:electronic-effects:9}
Wijs het \omterm{def:b1:electronic-effects:nucleophile}{nucleofiel} en het \omterm{def:b1:electronic-effects:nucleophile}{elektrofiel} aan in: \ce{CH3Br + OH-};
\ce{H2O + H+}; ethanal met een cyanide-ion; \ce{BF3 + NH3}. Teken de
\omterm{def:b1:electronic-effects:curly-arrow}{gebogen pijlen}.
\end{exercise}

\begin{exercise}[$\star\star\star$]\label{exo:b1:electronic-effects:10}
Een oplossing bevat \qty{0.010}{mol/L} trichloorethaanzuur. Is de formule
voor een \omterm{def:b1:predominant-reaction:strong-acid}{zwak zuur}, $\mathrm{pH} = \frac12(\mathrm{p}K_a - \log C)$,
geldig? Bereken de pH correct en de \omterm{def:b1:predominant-reaction:dissociation}{dissociatiegraad}.
\end{exercise}

\begin{exercise}[$\star\star\star$]\label{exo:b1:electronic-effects:11}
De $\mathrm{p}K_a$-waarden van alcoholen, gemeten in water, stijgen van
methanol (15.5) over ethanol (15.9) en propaan-2-ol (17.1) tot
2-methylpropaan-2-ol (19.2). Geef twee verklaringen, één via het
\omterm{def:b1:electronic-effects:inductive}{inductieve effect} van de alkylgroepen op het alkoxide, één via de
\omterm{def:b1:intermolecular-forces-solvents:dissolution}{solvatatie} van het alkoxide.
% ledger: pka:methanol, pka:ethanol, pka:2-propanol, pka:tBuOH
\end{exercise}


\begin{exercise}[$\star\star\star$]\label{exo:b1:electronic-effects:12}
Leg met de \omterm{def:b1:lewis-resonance-vsepr:resonance}{resonantiestructuur} van het amide uit waarom amiden op
stikstof basisch noch \omterm{def:b1:electronic-effects:nucleophile}{nucleofiel} zijn, terwijl aminen beide zijn. Welk
atoom van een amide wordt in \omterm{def:b1:predominant-reaction:strong-acid}{sterk zuur} geprotoneerd?
\end{exercise}

\section{Probleem: waarom trifluorethaanzuur zo sterk is}

\begin{problem}[{Weekendprobleem --- de inductieve reeks van de
gehalogeneerde zuren, carboxylaten tegenover alkoxiden, de drie
nitrofenolen, een berekening met de predominante reactie, en de
verhouding van de zuurconstanten van trifluorethaanzuur en ethaanzuur}]
\label{pb:b1:electronic-effects:1}
Gegevens ($\mathrm{p}K_a$ bij \qty{25}{\celsius}): ethaanzuur 4.76,
chloorethaanzuur 2.87, dichloorethaanzuur 1.35, trichloorethaanzuur 0.51,
trifluorethaanzuur 0.52, ethanol 15.9, fenol 9.99, 2-nitrofenol 7.23,
3-nitrofenol 8.36, 4-nitrofenol 7.16; $\mathrm{p}K_e = 14.00$.

\medskip
\noindent\textbf{Deel I --- De inductieve reeks.}
\begin{enumerate}
  \item Bereken de verhouding van de $K_a$ voor elke stap van ethaanzuur
        tot trichloorethaanzuur.
  \item Is het effect van elk chlooratoom hetzelfde? Becommentarieer.
  \item Verklaar de trend met het \omterm{def:b1:electronic-effects:inductive}{inductieve effect}.
  \item Waarom zou een chlooratoom op het koolstofatoom van een langere
        keten dat het verst van COOH ligt, bijna geen effect hebben?
  \item Fluor is elektronegatiever dan chloor. Wat toont de vergelijking
        van het trifluor- en het trichloorzuur, en waarom is het moeilijk
        zulke $\mathrm{p}K_a$-waarden nauwkeurig te meten?
  \item In welke vorm komt elk van deze zuren voor bij pH 3?
\end{enumerate}

\medskip
\noindent\textbf{Deel II --- Carboxylaten en alkoxiden.}
\begin{enumerate}[resume]
  \item Teken de twee \omterm{def:b1:lewis-resonance-vsepr:resonance}{resonantiestructuren} van het ethanoaation.
  \item Wat voorspellen ze voor de lengten van de twee C--O-bindingen van
        het ion?
  \item Waarom heeft het ethoxide-ion geen gelijkwaardige structuren?
  \item Bereken de verhouding van de $K_a$ van ethaanzuur en ethanol.
  \item Plaats fenol tussen beide en verantwoord met de
        \omterm{def:b1:lewis-resonance-vsepr:resonance}{resonantiestructuren} van het fenoxide.
  \item Welke van ethanol, fenol en ethaanzuur reageren bijna volledig met
        een oplossing van waterstofcarbonaat ($\mathrm{p}K_a$ 6.37 voor
        \ce{CO2,H2O}/\ce{HCO3-})?
\end{enumerate}

\medskip
\noindent\textbf{Deel III --- De nitrofenolen.}
\begin{enumerate}[resume]
  \item Rangschik de drie nitrofenolen en fenol naar zuurgraad.
  \item Teken een \omterm{def:b1:lewis-resonance-vsepr:resonance}{resonantiestructuur} van 4-nitrofenoxide met de lading op
        de nitrogroep.
  \item Toon aan dat zo'n structuur voor 3-nitrofenoxide niet bestaat.
  \item Waarom is 3-nitrofenol toch zuurder dan fenol?
  \item 2-Nitrofenol is iets zwakker dan 4-nitrofenol, hoewel beide het
        \omterm{def:b1:electronic-effects:mesomeric}{mesomere effect} hebben; stel een verklaring voor waarbij zijn OH
        en de nabije nitrogroep betrokken zijn.
  \item Welke nitrofenolen zijn bij pH 7.5 vooral in de anionvorm? Waarom
        maakt dat van 4-nitrofenol een \omterm{def:b1:titration-methods:indicator}{kleurindicator}?
\end{enumerate}

\medskip
\noindent\textbf{Deel IV --- Trifluorethaanzuur in water.}
\begin{enumerate}[resume]
  \item Bereken de pH van trifluorethaanzuur van \qty{0.010}{mol/L} (neem
        geen zwakke dissociatie aan).
  \item Bereken de pH van ethaanzuur van \qty{0.010}{mol/L}.
  \item Schrijf de reactie van trifluorethaanzuur met natriumethanoaat op
        en bereken haar constante.
  \item Geef de verhouding $K_a(\ce{CF3COOH})/K_a(\ce{CH3COOH})$, als macht
        van tien en als getal.
\end{enumerate}
\end{problem}
